% Foundations IV — shared manuscript macros.
% Stable, cross-section macros. Add macros incrementally as new
% vocabulary stabilizes; never redefine locally inside a section.
% Governed by the shared notation and terminology discipline.

%--- axis labels (the 8 primitive-necessity axes of the F-IV catalog) -------
\newcommand{\axisAccess}{Access}
\newcommand{\axisClosure}{Closure}
\newcommand{\axisStability}{Stability}
\newcommand{\axisTransport}{Transport}
\newcommand{\axisStatus}{Status}
\newcommand{\axisRecords}{Records}
\newcommand{\axisSymmetry}{Symmetry}
\newcommand{\axisInterfaces}{Interfaces}
\newcommand{\axisSelfReference}{Self-reference}

%--- chapter groupings (the 6 drafting chapters that group the 8 axes) ------
\newcommand{\chapAccess}{Access}
\newcommand{\chapStability}{Stability}
\newcommand{\chapTransport}{Transport}
\newcommand{\chapSymmetry}{Symmetry}
\newcommand{\chapStatusRecordsCoherence}{Status, Records, Coherence}
\newcommand{\chapSelfReferenceLimits}{Self-Reference and Limits}
\newcommand{\chapClosureRealityBoundary}{Closure--Reality Boundary}

%--- catalog identity --------------------------------------------------------
\newcommand{\FIV}{Foundations~IV}
\newcommand{\FoundI}{Foundations~I}
\newcommand{\FoundII}{Foundations~II}
\newcommand{\FoundIII}{Foundations~III}
\newcommand{\NumLaws}{52}
\newcommand{\NumMechanizeNow}{46}
\newcommand{\NumAnchor}{2}
\newcommand{\NumPartial}{2}
\newcommand{\NumObligation}{2}
\newcommand{\NumAxes}{8}
\newcommand{\NumChapters}{6}

%--- quotient discipline (shared Transport / Access / Stability) ------------
\newcommand{\currQ}{\pi^{\mathrm{cur}}}  % current quotient map
\newcommand{\predM}{\pi^{\mathrm{pred}}} % predictive quotient map
\newcommand{\splitObs}{\Delta}         % split-pair obstruction
\newcommand{\quotMap}{\pi}             % quotient map
\newcommand{\accessQ}{q_I}             % access quotient
\newcommand{\obsSig}{\sigma_I}         % observation signature
\newcommand{\closureForm}{\mathsf{cl}} % formed-closure operator
\newcommand{\BirdInt}{\mathsf{BirdInt}}% interaction judgment

%--- substrate-anchor notation (Papers 1, 3 -- present in inherited Lean) ---
% Paper 1 (Adequacy Residuals) — for F6/F8 anchor paragraphs.
% Note: \Xi is a built-in Greek letter; we use it directly without
% redefinition. The adequacy residual is written $\Xi(D \mid L)$.
\newcommand{\Kdom}{K^{D}}              % dominator
\newcommand{\Kloss}{K^{L}}             % loss
\newcommand{\Schur}{\mathsf{Schur}}    % Schur identity
% Paper 3 (AOR) — for F53 anchor paragraph
\newcommand{\ClSB}{\mathrm{Cl}_{\mathrm{SB}}}  % reflective closure operator
\newcommand{\etaSB}{\eta_{\mathrm{SB}}} % AOR unit
\newcommand{\XiSB}{\Xi_{\mathrm{SB}}}   % AOR residual

%--- substrate-anchor notation (Papers 2, 6 -- substrate_pending in Lean) ----
% Paper 2 (SDTC) / Paper 6 (RH) — for F14 paragraph
\newcommand{\Jinv}{J}                  % involution
\newcommand{\Fixinv}{\mathrm{Fix}(\Jinv)}  % involution fixity set
\newcommand{\trAX}{\operatorname{tr}}  % trace functional

%--- normal-form template (uniform across all 52 laws) ----------------------
% Use these environments to wrap the catalog laws in body prose.
% Each F-IV law follows the 6-part template: Setup → Theorem → Proof
% spine → Case enumeration → BirdInt status → Nonclaims.
\newcommand{\NormalForm}[1]{\textbf{Normal form #1}}
\newcommand{\SetupBlock}{\textit{Setup.}}
\newcommand{\ProofSpine}{\textit{Proof spine.}}
\newcommand{\CaseEnum}{\textit{Cases.}}
\newcommand{\BirdIntStatus}{\textit{Interaction status.}}
\newcommand{\Nonclaims}{\textit{Nonclaims.}}

%--- status tags (display form for paper text) ------------------------------
% Body prose uses these textual descriptions; the literal Lean tag
% strings (`# obligation`, `# substrate_pending`) appear only in
% app:formalization per feedback_academic_register.md.
\newcommand{\statTheorem}{Theorem}
\newcommand{\statPartial}{Partial}
\newcommand{\statObligation}{obligation (axiom, opacity-guarded)}
\newcommand{\statAnchor}{Anchor}
\newcommand{\statSubstratePending}{substrate-pending (axiom-applied)}

%--- anchor cross-reference shorthands --------------------------------------
% These are paper-prose hooks for the substrate papers. They expand
% to descriptive title phrases, not the in-house ``Paper N'' series
% numbering. Bibtex keys live in paper/references.bib.
\newcommand{\PaperAR}{the adequacy-residuals paper}
\newcommand{\PaperSDTC}{the self-dual trace confinement paper}
\newcommand{\PaperAOR}{the audited-operational-realisability paper}
\newcommand{\PaperNeedleKiller}{the needle-killer paper}
\newcommand{\PaperNS}{the Navier--Stokes paper}
\newcommand{\PaperRH}{the Riemann Hypothesis paper}
\newcommand{\PaperPvNP}{the P versus NP paper}
\newcommand{\PaperCSL}{the CSL paper}

%--- opacity-guard predicate names (display form in body prose) -------------
% In body text the predicate is named descriptively; the literal Lean
% predicate name appears only in app:formalization.
\newcommand{\GuardPvNP}{saturated-SAT source predicate}
\newcommand{\GuardCSL}{CSL diagnosis source predicate}
\newcommand{\GuardDuality}{duality-fixity substrate-anchor predicate}
\newcommand{\GuardSSB}{future branch-selection formation source predicate}

%--- canonical theorem-name shortcuts (for cross-references in body) --------
% These render the canonical paper-prose names so cross-references are
% mechanically enforced. Add as needed during drafting.
\newcommand{\thmDescentRepair}{Descent--Repair Normal Form}
\newcommand{\thmHolonomyMemory}{Holonomy-Memory Repair Normal Form}
\newcommand{\thmLocalGlobalObstruction}{Local--Global Obstruction Normal Form}
\newcommand{\thmRecombinationComparison}{Recombination Comparison}
\newcommand{\thmNoNeedles}{No-Needles Stability}
\newcommand{\thmSufficiencyClosure}{Sufficiency Closure}
\newcommand{\thmStrictExtension}{Strict Extension}
\newcommand{\thmNoUnstatusedResidual}{No Unstatused Residual}
\newcommand{\thmQuotientality}{Quotientality}
\newcommand{\thmAdequacy}{Adequacy / No-Overread}
\newcommand{\thmNoFreeDistinction}{No-Free-Distinction}
\newcommand{\thmHiddenness}{Hiddenness Normal Form}
\newcommand{\thmCSLFormedClosure}{CSL Formed-Closure Law}
\newcommand{\thmLayerMultiplicity}{Layer Multiplicity}
\newcommand{\thmDualityFixity}{Duality Fixity}
\newcommand{\thmSelectionObstruction}{Selection Obstruction}
\newcommand{\thmSpontaneousSSB}{Spontaneous Symmetry Breaking}
\newcommand{\thmAnomalySymmetry}{Anomaly / Symmetry Obstruction}
\newcommand{\thmDualityEquivalence}{Duality Equivalence}
\newcommand{\thmFactRecordFormation}{Fact / Record Formation}
\newcommand{\thmHorizonSufficiency}{Horizon Sufficiency}
\newcommand{\thmClosureRealityBoundary}{Closure--Reality Boundary}
\newcommand{\thmClosureCompletenessBoundary}{Closure Completeness Boundary}
% Additional thm* shortcuts are added only when body prose needs them.

% F-code of a catalog row, printed as a link to its numbered result.
\makeatletter
\newcommand{\Fref}[1]{\@ifundefined{r@thm:#1}{\@ifundefined{r@lem:#1}{\@ifundefined{r@prop:#1}{#1}{\hyperref[prop:#1]{#1}}}{\hyperref[lem:#1]{#1}}}{\hyperref[thm:#1]{#1}}}
\makeatother
